\documentclass[11pt,a4paper]{article}
\usepackage[utf8]{inputenc}
\usepackage[T1]{fontenc}
\usepackage{geometry}
\geometry{margin=2.3cm}

\title{Rapport d'avancement\\GNU -- Gestion des notes universitaires}
\author{Projet GNU -- Département d'Informatique}
\date{13 octobre 2026}

\begin{document}
\maketitle
\tableofcontents
\bigskip

\section{Résumé exécutif}
Le projet dispose maintenant d'une base PostgreSQL structurée et alimentée par les programmes harmonisés Licence/Master, ainsi que d'un frontend Flutter démontrable pour les trois rôles : étudiant, enseignant et agent informatique. Les données de la classe INF361--SR ont été importées pour l'année 2025--2026 et les inscriptions ont été validées.

Les principaux travaux restants concernent l'intégration réelle entre Flutter, l'API Laravel et PostgreSQL, l'authentification, la gestion des autorisations et la couverture complète des écrans métier.

\section{Travaux réalisés}
\subsection{Sauvegarde et organisation}
\begin{itemize}
  \item Ancienne version conservée dans le dossier de sauvegarde.
  \item Données organisées par année académique.
  \item Années de référence enregistrées : 2021--2022, 2024--2025 et 2025--2026.
\end{itemize}

\subsection{Modèle académique}
\begin{itemize}
  \item Filière canonique : \textbf{INF -- Informatique}.
  \item Spécialités définies : \textbf{GL}, \textbf{SR}, \textbf{SECU} et \textbf{DATA}.
  \item L1 et L2 considérés comme tronc commun ; la spécialité est portée par le niveau ou la classe lorsque la source le précise.
  \item L'ancienne référence Génie logiciel a été conservée comme historique.
  \item INF331 est traitée comme UE fondamentale de la filière Informatique et non comme une spécialité.
\end{itemize}

\subsection{Catalogue UE et EC}
\begin{itemize}
  \item 73 UE importées dans le catalogue Informatique.
  \item 121 EC extraits des programmes harmonisés et importés.
  \item 62 UE activées après contrôle de cardinalité.
  \item Règle appliquée : une UE active possède un ou deux EC.
  \item Les crédits sont distingués des coefficients ; aucun coefficient artificiel n'a été déduit des PDF.
  \item Les codes EC du Master sont des identifiants techniques dérivés du code UE et du rang, car les grilles ne fournissent pas de codes EC officiels.
\end{itemize}

\subsection{Année 2025--2026 et étudiants}
\begin{itemize}
  \item Classes Informatique créées pour L1, L2, L3, M1 et M2.
  \item Classe opérationnelle : \texttt{INF-L3-2025-2026}.
  \item 35 étudiants INF361--SR importés.
  \item 35 inscriptions en classe créées et validées.
  \item 350 inscriptions UE créées et validées.
  \item Les données historiques de \texttt{L3-GL-A} sont conservées conformément à la décision de ne pas supprimer les références protégées.
\end{itemize}

\subsection{Frontend Flutter}
\begin{itemize}
  \item Parcours de démonstration pour étudiant, enseignant et agent informatique.
  \item Navigation responsive testée sur petit et grand écran.
  \item Modules de démonstration : tableau de bord, résultats, inscription, notes, évaluations, requêtes, rattrapages, programmes, comptes, bulletins, audit et publication.
  \item Thème clair/sombre ajouté avec bascule depuis la barre supérieure.
  \item Corrections des assertions Flutter liées aux composants \texttt{ListTile} et aux surfaces Material.
  \item Tests Flutter validés : 13 tests réussis.
  \item Build Web validé avec \texttt{flutter build web --no-pub}.
\end{itemize}

\section{Scripts et documents produits}
\begin{tabular}{p{0.43\textwidth}p{0.48\textwidth}}
Fichier & Rôle\\\hline
001\_structure\_academique.sql & structure université, filière, niveaux et années\\
002\_import\_ue\_catalogue.sql & import idempotent du catalogue UE\\
003\_specialites\_informatique.sql & spécialités GL, SR, SECU et DATA\\
005\_rattachement\_historique\_gl.sql & rattachement historique de l'ancienne classe\\
006\_import\_ec\_catalogue.sql & import et activation contrôlée des EC\\
007\_classes\_2025\_2026.sql & création des classes annuelles\\
009\_import\_etudiants\_inf361\_2025\_2026.sql & import étudiants et inscriptions brouillon\\
010\_valider\_inscriptions\_inf\_l3\_2025\_2026.sql & validation administrative\\
REGLES\_CATALOGUE.md & décisions de normalisation et règles métier\\
\end{tabular}

\section{Travaux restant à faire}
\subsection{Priorité 1 -- API et sécurité}
\begin{itemize}
  \item Finaliser l'API Laravel et ses connexions PostgreSQL.
  \item Remplacer le \texttt{DemoStore} Flutter par des services HTTP réels.
  \item Implémenter connexion, inscription, renouvellement et réinitialisation de mot de passe.
  \item Appliquer les autorisations par rôle et par classe/matière.
  \item Ne plus utiliser de mots de passe temporaires générés lors de l'import étudiant.
\end{itemize}

\subsection{Priorité 2 -- Fonctionnalités métier}
\begin{itemize}
  \item Affectation des enseignants aux EC et classes.
  \item Configuration des évaluations et pondérations.
  \item Saisie, contrôle, validation et publication des notes.
  \item Calcul des résultats, bulletins semestriels et annuels.
  \item Requêtes, décisions, rattrapages et journal d'audit en données réelles.
\end{itemize}

\subsection{Priorité 3 -- Qualité et exploitation}
\begin{itemize}
  \item Corriger les avertissements de clic hors écran dans les tests UI avec \texttt{ensureVisible}.
  \item Ajouter tests d'intégration API, tests de sécurité et tests de non-régression.
  \item Documenter la procédure de démarrage PostgreSQL, Laravel et Flutter.
  \item Ajouter sauvegardes automatisées et stratégie de restauration.
  \item Vérifier les écarts de crédits Master actuellement documentés pour INF461, INF561 et INF571 avant publication.
\end{itemize}

\section{État global}
\begin{tabular}{p{0.58\textwidth}p{0.25\textwidth}}
\textbf{Domaine} & \textbf{État}\\\hline
Structure académique & Réalisé\\
Catalogue UE/EC & Réalisé avec contrôles documentaires restants\\
Classes et inscriptions 2025--2026 & Réalisé pour INF361--SR\\
Frontend de démonstration & Réalisé et testé\\
Thème clair/sombre & Réalisé\\
API Laravel réelle & À faire\\
Authentification et sécurité & À faire\\
Notes, bulletins et publication réels & À faire\\
\end{tabular}

\section{Conclusion}
La fondation fonctionnelle et documentaire du projet est en place. Le prochain jalon est l'intégration API : le frontend devra lire les niveaux, UE, EC, classes, inscriptions et rôles depuis PostgreSQL via Laravel, puis appliquer les mêmes règles de validation que les scripts d'import.

\end{document}
